\documentclass{article}
\usepackage[english,greek]{babel}
\usepackage[T1]{fontenc}
\usepackage[utf8]{inputenc}
\usepackage{amsmath}

\title{\selectlanguage{greek}Eisagwg'h sthn analutik'h jewr'ia ariqm'wn}
\author{\selectlanguage{greek}N. Papadopoulos}
\date{2026}

\begin{document}
\maketitle

\selectlanguage{greek}

\section{Eisagwg'h}

H analutik'h jewr'ia ariqm'wn sqet'izetai me thn katanom'h twn
pr'wtwn ariqm'wn.  To k'urio erg'aleio e'inai h sun'arthsh z'hta
tou Riemann, pou or'izetai ws

\selectlanguage{english}
\[
  \zeta(s) = \sum_{n=1}^{\infty} \frac{1}{n^s}, \quad \text{Re}(s) > 1.
\]
\selectlanguage{greek}

% Greek text with oxia/tonos distinction
To pr'oblhma e'inai na deiqje'i 'oti h sun'arthsh $\zeta(s)$
den mhden'izetai sthn euqe'ia $\text{Re}(s) = 1$.

\section{To Je'wrhma twn Pr'wtwn Ariqm'wn}

\selectlanguage{english}
\begin{theorem}
Let $\pi(x)$ denote the number of primes up to~$x$.  Then
$\pi(x) \sim x / \ln x$ as $x \to \infty$.
\end{theorem}
\selectlanguage{greek}

To je'wrhma apode'iqjhke anex'arthta ap'o ton Hadamard kai ton
de la Vall'ee Poussin to 1896.  H ap'odeixh qrhsimopoi'ei tis
analutik'es idi'othtes ths $\zeta(s)$.

\section{Sum'erasma}

H up'oqesh tou Riemann param'enei 'ena ap'o ta pio shmantiok'a
an'oluita probl'hmata twn majhmatik'wn.

\end{document}
